\documentclass[11pt]{article}

\usepackage[T1]{fontenc}
\usepackage[utf8]{inputenc}
\usepackage{lmodern}
\usepackage{amsmath,amssymb,amsthm,mathtools}
\usepackage{geometry}
\usepackage{hyperref}
\usepackage{enumitem}

\geometry{margin=1in}
\setlist[itemize]{leftmargin=1.5em}
\setlist[enumerate]{leftmargin=1.75em}
\emergencystretch=3em

\newcommand{\ec}[1]{\texttt{\detokenize{#1}}}
\newcommand{\Fq}{\mathbf F_q}
\newcommand{\Z}{\mathbb Z}
\newcommand{\br}{\operatorname{br}_8}
\newcommand{\zetaNTT}{\zeta}

\newtheorem{theorem}{Theorem}
\newtheorem{lemma}{Lemma}
\newtheorem{proposition}{Proposition}
\newtheorem{obligation}{Obligation}
\newtheorem{remark}{Remark}

\title{End-to-End HAETAE NTT Correctness: Checked Core and Remaining Algebraic Bridge}
\author{Codex Proof Notes}
\date{29 April 2026}

\begin{document}
\maketitle

\begin{abstract}
This note states the intended end-to-end correctness theorem for the HAETAE
Jasmin NTT, records the portions that are currently machine-checked in
EasyCrypt, and identifies the exact missing algebraic bridge.  The main point is
that the already checked pRHL development proves the extracted Jasmin cores
against the imperative procedures \ec{NTT_Fq.NTT.ntt} and
\ec{NTT_Fq.NTT.invntt}; it does not yet prove those imperative procedures
against the full closed-form polynomial NTT.  The file \ec{NTTFullSpec.ec}
therefore introduces the correct full transform and proves that the older
\ec{Rq.ntt} operator is not the full HAETAE NTT.
\end{abstract}

\section{The Correct Abstract Forward Transform}

Let \(q=64513\), let \(\Fq=\Z/q\Z\), and let
\(\zetaNTT=\ec{zroot}=\ec{incoeff 426}\).  For a coefficient array
\(p\in\Fq^{256}\), the full HAETAE negacyclic NTT is the array
\[
  \widehat p_i
  =
  \sum_{j=0}^{255}
    p_j\,\zetaNTT^{(2\br(i)+1)j},
  \qquad 0\le i<256,
\]
where \(\br(i)=\ec{bsrev 8 i}\).  This is formalized in
\ec{NTTFullSpec.full_ntt}.

\begin{lemma}[bit-reversed partial specification]
The predicate \ec{NTTFullSpec.full_ntt_spec r p} states equivalently that
\[
  r_{\br(s)}
  =
  \sum_{j=0}^{255}
    p_{\br(j)}\,
    \zetaNTT^{(2s+1)\br(j)}
  \qquad (0\le s<256).
\]
\end{lemma}

\begin{proof}
This is the definition of \ec{partial_ntt} at length \(256\), block index
\(0\), and start index \(s\).  The summation is indexed in bit-reversed input
order because the loop-level NTT proof naturally tracks the array positions
visited by the Cooley--Tukey butterflies.
\end{proof}

\begin{proposition}[checked in \ec{NTTFullSpec.ec}]
\[
  \ec{full_ntt_spec r p}\quad\Longrightarrow\quad
  r=\ec{full_ntt p}.
\]
\end{proposition}

\begin{proof}
The EasyCrypt lemma is \ec{full_ntt_spec_imp}.  For an output index \(i\), apply
the specification at \(s=\br(i)\).  Since bit reversal is an involution on
\(\{0,\ldots,255\}\), the left-hand side becomes \(r_i\).  The right-hand side
is then reindexed along the permutation \(j\mapsto \br(j)\).  Commutativity of
the finite sum and the identity \(\br(\br(j))=j\) convert the bit-reversed
partial sum into the closed-form coefficient of \ec{full_ntt}.
\end{proof}

\section{Why the Existing \texorpdfstring{\ec{Rq.ntt}}{Rq.ntt} Is Not Enough}

The operator \ec{Rq.ntt} is an incomplete even/odd transform:
even output coefficients sum only even input coefficients, and odd output
coefficients sum only odd input coefficients.  In formulas,
\[
  (\ec{Rq.ntt}(p))_{2i}
  =
  \sum_{j=0}^{127}
    p_{2j}\,\zetaNTT^{(2\br(i)+1)j},
  \qquad
  (\ec{Rq.ntt}(p))_{2i+1}
  =
  \sum_{j=0}^{127}
    p_{2j+1}\,\zetaNTT^{(2\br(i)+1)j}.
\]
This is not the full \(256\)-point transform above.

\begin{proposition}[checked in \ec{NTTFullSpec.ec}]
\[
  \ec{Rq.ntt Rq.one} \ne \ec{NTTFullSpec.full_ntt Rq.one}.
\]
\end{proposition}

\begin{proof}
At output index \(1\), the incomplete transform sees only odd input
coefficients.  The polynomial \(\ec{Rq.one}\) has coefficient \(1\) at index
\(0\) and \(0\) elsewhere, so \((\ec{Rq.ntt Rq.one})_1=0\).  The full transform
at index \(1\), however, contains the \(j=0\) term \(1\cdot \zetaNTT^0=1\) and
all other terms vanish, so \((\ec{full_ntt Rq.one})_1=1\).  Since \(1\ne0\) in
\(\Fq\), the arrays are unequal.  This is the EasyCrypt lemma
\ec{rq_ntt_is_not_full_ntt_on_one}.
\end{proof}

\section{Checked Implementation Refinement Layer}

The file \ec{RefJasminNTT.ec} proves the direct pRHL connection from the
extracted Jasmin cores to the imperative EasyCrypt model.

\begin{theorem}[forward core refinement, checked]
If the input byte array \(\rho\) represents \(p\) with 16-bit word bounds, then
running \ec{Hpoly_extract.M._poly_ntt} on \(\rho\) and running
\ec{NTT_Fq.NTT.ntt} on \(p\) with \ec{NTT_Fq.zetas} produce related outputs:
\[
  \ec{poly_repr_bound}(\rho,p,16)
  \Longrightarrow
  \ec{poly_repr_bound}(\rho',r',24).
\]
\end{theorem}

\begin{proof}
This is \ec{RefJasminNTT.poly_ntt_core_ref}.  The proof synchronizes the three
nested loops of the extracted Jasmin core with the three nested loops of
\ec{NTT_Fq.NTT.ntt}.  The invariant tracks equal loop counters, equality of
twiddle indices, the table bridge between Montgomery-encoded machine constants
and abstract twiddles, and pointwise bounds showing that word addition,
subtraction, and Montgomery multiplication do not overflow.  Each butterfly is
closed by a local field equality:
\[
  \operatorname{word}(\ec{fqmul}(z,x))
  =
  \operatorname{word}(z)\operatorname{word}(x)R^{-1},
\]
followed by cancellation of \(R\) using the table lemmas.
\end{proof}

\begin{theorem}[inverse core refinement, checked]
If the input byte array \(\rho\) represents \(p\) with 16-bit word bounds, then
\ec{Hpoly_extract.M._poly_invntt} refines \ec{NTT_Fq.NTT.invntt}; the concrete
output represents the Montgomery-scaled version of the imperative inverse
output:
\[
  \ec{poly_repr_bound}(\rho,p,16)
  \Longrightarrow
  \ec{poly_repr_bound}(\rho',\ec{array256_mont}(r'),16).
\]
\end{theorem}

\begin{proof}
This is \ec{RefJasminNTT.poly_invntt_core_ref}.  The butterfly-loop argument is
the inverse analogue of the forward proof.  The final scaling loop is handled by
a separate prefix invariant: processed coefficients have been multiplied by the
special table entry at index \(255\), while unprocessed coefficients still
represent the pre-scaling abstract state.  The special table lemma
\[
  \operatorname{word}(\ec{jzetas_inv}[255])R^{-1}
  =
  \ec{scale255}\,R
\]
explains the Montgomery factor in the postcondition.
\end{proof}

\section{The Missing End-to-End Bridge}

The inverse side now has the same specification-side conversion as the forward
side.  The EasyCrypt predicate \ec{full_invntt_spec r p} states the inverse
closed form at bit-reversed output indices, and the checked lemma
\ec{full_invntt_spec_imp} proves that this predicate implies
\(r=\ec{full_invntt p}\).  Thus the remaining inverse work is not a
closed-form reindexing problem; it is the loop algebra showing that
\ec{NTT_Fq.NTT.invntt} establishes \ec{full_invntt_spec}.

To obtain the full forward theorem for the Jasmin core, one still needs the
following Hoare statement.

\begin{obligation}[forward imperative algebra]
\[
  \ec{NTT_Fq.NTT.ntt}(p,\ec{NTT_Fq.zetas})
  =
  \ec{NTTFullSpec.full_ntt}(p).
\]
\end{obligation}

Equivalently, a loop proof may establish \ec{full_ntt_spec res p} at the end of
\ec{NTT_Fq.NTT.ntt}, then use \ec{full_ntt_spec_imp}.  The stage invariant must
use the full-transform exponent
\[
  \br(s\bmod \ell),
\]
not the MLKEM incomplete-transform exponent
\[
  \br(2(s\bmod \ell)).
\]
The latter only describes the even/odd split transform.

\begin{obligation}[inverse imperative algebra]
\[
  \ec{NTT_Fq.NTT.invntt}(p,\ec{NTT_Fq.zetas_inv})
  =
  \ec{NTTFullSpec.full_invntt}(p).
\]
\end{obligation}

This inverse theorem must also be reconciled with the Montgomery-scaled
postcondition of \ec{poly_invntt_core_ref}.  The composed Jasmin theorem should
therefore either state the output relation against
\(\ec{array256_mont}(\ec{full_invntt}(p))\), or use a representation relation
whose codomain is explicitly Montgomery encoded.

\section{Conclusion}

The present EasyCrypt development is strong enough to justify the statement:

\begin{quote}
The extracted Jasmin HAETAE NTT cores have been machine-checked against the
imperative EasyCrypt NTT procedures, with word-level overflow obligations and
Montgomery table bridges included.
\end{quote}

It is not yet strong enough to justify the stronger statement:

\begin{quote}
The complete HAETAE Jasmin NTT implementation has been machine-checked against
the full abstract polynomial NTT specification.
\end{quote}

That stronger claim becomes valid only after the two imperative algebra
obligations above are proved in EasyCrypt without \ec{admit}, \ec{abort}, or new
axioms and then composed with \ec{poly_ntt_core_ref} and
\ec{poly_invntt_core_ref}.

\end{document}
